% ЗГЕНЕРОВАНО scripts/build_format_v2_chapter.py — не редагувати вручну.
% Джерело: editorial/cards/12.*.json, format-v2/rozdil-12.json, paraleli.json.

\naukchapter{12}{8}{Глава 12. Саме в цьому я стою}
\naukintro{Дванадцята глава вся закінчується тими самими словами — саме в цьому я стою, — і цей рефрен тут не приспів, а мірило: кожен вірш спершу перелічує, чого людина вже не робить, і лише потім каже, де вона стоїть. Перші три вірші йдуть від тіла до мови й далі до думки; наступні три знімають пари — приймати й відкидати, обряд і його припинення. Останній вірш не додає нічого нового: він просто називає того, хто так стоїть.}
\vsec{12.1}
\vdeva{कायकृत्यासहः पूर्वं ततो वाग्विस्तरासहः । अथ चिन्तासहस्तस्माद् एवम् एवाहम् आस्थितः ॥}
\vtranslit{каякрітясахах пурвам тато ваґвістарасахах | атха чінтасахастасмад евам евахам астхітах ||}
\vperek{Спершу нездатний до тілесних справ, потім нездатний до поширення мови, далі нездатний до думання — тому саме в цьому я стою.}
\vcom{Тут відмова йде **не за бажанням, а через неспроможність**: `asahaḥ` — «нездатний витримати», а не «той, хто відмовився». Це рідкісний для тексту випадок, коли стан описано знизу — через те, чого людина вже не може. Формула `evam evāham āsthitaḥ` повторюється в 12.1–12.4, як `niścayī` у главі 11: переклад мусить тримати її однаково.}
\vpar{Українська паралель (наша інтерпретація)}{Василь Стус, «Час творчості».}{Наша паралель: у Стуса нездатність до звичних справ теж постає не як втрата, а як те, у чому людина нарешті стоїть. Наше порівняння — типологічне.}
\vfig{ilustratsii/12-1-neti.pdf}{Нездатність до тіла, мови й думки — і стояння в цьому}{Схема наша, не історична ілюстрація.}

\vsec{12.2}
\vdeva{प्रीत्यभावेन शब्दादेर् अदृश्यत्वेन चात्मनः । विक्षेपैकाग्र्यहृदय एवम् एवाहम् आस्थितः ॥}
\vtranslit{прітябгавена шабдадер адрішятвена чатманах | вікшепайкаґр'яхрідая евам евахам астхітах ||}
\vperek{Через відсутність прихильності до звуку й подібного та через невидимість його для самого себе — серце, для якого розсіяність і зосередженість одне — саме в цьому я стою.}
\vcom{`vikṣepa-aikāgrya` в одному складнику — найважливіше місце вірша: розсіяність і зосередженість названі **одним** словом, і це не поетична вільність, а твердження, що для того, хто стоїть у собі, різниці між ними немає. `adṛśyatva` («невидимість») стосується `ātman` — себе не видно собі як об'єкт.}
\vpar{Українська паралель (наша інтерпретація)}{Леся Українка, «У пущі» (1909).}{Наша паралель: у Лесі Українки зосередженість і розсіяність теж виявляються одним станом, а не двома. Наше порівняння — за будовою думки.}
\vfig{ilustratsii/12-2-witness.pdf}{Серце в центрі: розсіяність і зосередженість — одне}{Схема наша, не історична ілюстрація.}

\vsec{12.3}
\vdeva{समाध्यासादिविक्षिप्तौ व्यवहारः समाधये । एवं विलोक्य नियमम् एवम् एवाहम् आस्थितः ॥}
\vtranslit{самадгясадівікшіптау в'явахарах самадгає | евам вілокя ніямам евам евахам астхітах ||}
\vperek{Коли розсіяність походить від прагнення самадгі та подібного, то й звичайна діяльність — задля самадгі. Так побачивши це правило, саме в цьому я стою.}
\vcom{Найдошкульніший вірш глави: **прагнення самадгі саме стає розсіяністю** (`vikṣipti`). І далі — іронія, яку легко проґавити: якщо навіть буденна діяльність робиться «задля самадгі», то правило (`niyama`) виявляється всеохопним, і виходу з нього немає. Переклад мусить зберегти цю замкненість, а не зробити з неї пораду.}
\vpar{Українська паралель (ПОТРЕБУЄ ПЕРЕВІРКИ)}{українська релігійна проза, не визначено.}{ПОТРЕБУЄ ПЕРЕВІРКИ: думка, що звичайна діяльність і самадгі — одне, близька до української релігійної прози, але конкретного твору ми не звірили.}
\vfig{ilustratsii/12-3-balance.pdf}{Дві шальки: прагнення самадгі й звичайна діяльність}{Схема наша, не історична ілюстрація.}

\vsec{12.4}
\vdeva{हेयोपादेयविरहाद् एवं हर्षविषादयोः । अभावाद् अद्य हे ब्रह्मन् नेवम् एवाहम् आस्थितः ॥}
\vtranslit{хейопадеявірахад евам харшавішадайох | абгавад адя хе брахман невам евахам астхітах ||}
\vperek{Через відсутність того, що відкидати й приймати, а отже — через відсутність радості й смутку, нині, о Брахмане, саме в цьому я стою.}
\vcom{`he brahman` — звертання до Джанаки як до брахмана, єдине таке в розділі. Ланцюг причинний: немає пари «відкидати/приймати» → немає радості й смутку → нема чого досягати. **Відкрите питання:** у SARIT наприкінці стоїть `nevamevāhamāsthitaḥ` (sandhi `na + evameva`), і я прочитав це як ствердне `evameva` за паралелізмом. Перевірка решти глави це підтверджує лише частково: 12.5, 12.6 і 12.7 усі закінчуються на `evamevāhamāsthitaḥ` **без** `na`, тобто 12.4 — єдиний виняток у главі. Отже `na` може бути або справжнім (і тоді сенс вірша протилежний), або помилкою транскрипції. **GRETIL недоступний, другого свідка немає** — картка лишається чернеткою, а не «перевірено».}
\vpar{Українська паралель (наша інтерпретація)}{Ліна Костенко, «Маруся Чурай» (1979).}{Наша паралель: у Костенко теж є стан, у якому немає що відкидати й немає чого приймати, і саме в ньому людина стоїть. Наше порівняння — типологічне.}
\vfig{ilustratsii/12-4-neti.pdf}{Відкидання й приймання, радість і смуток відкинуті}{Схема наша, не історична ілюстрація.}

\vsec{12.5}
\vdeva{आश्रमानाश्रमं ध्यानं चित्तस्वीकृतवर्जनम् । विकल्पं मम वीक्ष्य एतैर् एवम् एवाहम् आस्थितः ॥}
\vtranslit{ашраманашрамам дгянам чіттасвікрітаварджанам | вікалпам мама вікшя етайр евам евахам астхітах ||}
\vperek{Медитація зі стадією життя і без неї, прийняття й відкидання, зроблені умом, — усе це домислювання. Побачивши це в собі, через саме це я так стою.}
\vcom{Єдине місце в книзі, де `āśrama` (стадія життя) названо **поряд з її відсутністю** — і обидва варіанти відкинуті як `vikalpa`. Це важливо для нашого словника: текст не обирає між «життям у стадії» і «поза нею», він знімає саму пару. `citta-svīkṛta-varjana` — прийняття й відкидання не речей, а **вироблені умом**; переклад мусить зберегти, що йдеться про акт ума, а не про оцінку світу.}
\vpar{Українська паралель (ПОТРЕБУЄ ПЕРЕВІРКИ)}{українська філософська проза, не визначено.}{ПОТРЕБУЄ ПЕРЕВІРКИ: думка, що медитація й обряди є домислом, найближча до української філософської прози, але назвати твір без звірки ми не можемо.}
\vfig{ilustratsii/12-5-knot.pdf}{Вузол домислу: медитація й обряд}{Схема наша, не історична ілюстрація.}

\vsec{12.6}
\vdeva{कर्मानुष्ठानम् अज्ञानाद् यथैव उपरमस् तथा । बुद्ध्वा सम्यग् इदं तत्त्वम् एवम् एवाहम् आस्थितः ॥}
\vtranslit{кармануштханам аджнянад ятхайва упарамас татха | буддгва сам'яґ ідам таттвам евам евахам астхітах ||}
\vperek{Виконання обрядів через незнання і припинення їх — однаково. Зрозумівши правильно цю суть, я так стою.}
\vcom{`yathaiva … tathā` («так само») ставить дію й відмову від дії на один рівень — і обидві названі `ajñāna`, коли виходять із незнання. Це не заклик припинити обряди: текст каже, що припинення не краще за виконання. Переклад мусить звучати рівно, без інтонації «краще не робити».}
\vpar{Українська паралель (ПОТРЕБУЄ ПЕРЕВІРКИ)}{українське бароко, не визначено.}{ПОТРЕБУЄ ПЕРЕВІРКИ: рівність обряду й припинення обряду — думка, близька до барокової полеміки, але конкретного тексту ми не перевірили.}
\vfig{ilustratsii/12-6-balance.pdf}{Дві шальки: обряд і його припинення — однаково}{Схема наша, не історична ілюстрація.}

\vsec{12.7}
\vdeva{अचिन्त्यं चिन्त्यमानो ऽपि चिन्तारूपं भजत्य् असौ । त्यक्त्वा तद्भावनं तस्माद् एवम् एवाहम् आस्थितः ॥}
\vtranslit{ачінтям чінтямано ’пі чінтарупам бгаджатй асау | тяктва тадбгаванам тасмад евам евахам астхітах ||}
\vperek{Незбагненне, навіть коли його осмислюють, набирає форми думки. Тому, відкинувши те плекання, я так стою.}
\vcom{Механізм названий точно: незбагненне не залишається собою, коли його осмислюють — воно **набирає форми думки** (`cintārūpa`). Тому відкидається не предмет, а `bhāvana` — плекання, повторюване обертання. Це та сама `cintā`, що в 11.5, і словник має бути спільним.}
\vpar{Українська паралель (наша інтерпретація)}{Григорій Сковорода, «Алфавіт, чи Буквар миру» (діалог).}{Наша паралель: у Сковороди незбагненне теж набирає форми думки, і саме тому його не можна втримати плеканням. Наше порівняння — типологічне.}
\vfig{ilustratsii/12-7-illusion.pdf}{Незбагненне під серпанком набирає форми думки}{Схема наша, не історична ілюстрація.}

\vsec{12.8}
\vdeva{एवम् एव कृतं येन स कृतार्थो भवेद् असौ । एवम् एव स्वभावो यः स कृतार्थो भवेद् असौ ॥}
\vtranslit{евам ева крітам єна са крітартхо бгавед асау | евам ева свабгаво ях са крітартхо бгавед асау ||}
\vperek{У кого це так зроблено — той стає тим, чию мету здійснено; у кого така природа — той стає тим, чию мету здійснено.}
\vcom{Вірш повторює рядок двічі й міняє одне слово: `kṛtam` (зроблено) → `svabhāvaḥ` (природа). Це не тавтологія, а розрізнення: спершу те, що людина зробила, потім те, чим вона стала. `kṛtārtha` — «той, чию мету здійснено», усталене слово; переклад тримає його однаково в обох половинах. Глава закінчується не настановою, а констатацією.}
\vpar{Українська паралель (наша інтерпретація)}{Микола Зеров, «Камена» (1924).}{Наша паралель: у Зерова завершеність теж постає як наслідок способу бачити, а не як наслідок зусилля. Наше порівняння — за будовою думки.}
\vfig{ilustratsii/12-8-flame.pdf}{Полумʼя здійсненої мети}{Схема наша, не історична ілюстрація.}

